% 均值漂移一维教学算例；宽约90mm。五个像素位于0至4，颜色灰灰灰红红；目标模型红0.6、灰0.4；y0=2，h=2.5，Epanechnikov轮廓。
\begin{tikzpicture}[x=1.45cm,y=1cm,font=\figurefont\small,
  arr/.style={-{Stealth[length=2mm]},BookInk,line width=0.9pt}]
  % 像素
  \foreach \x/\col in {0/BookMuted!30,1/BookMuted!30,2/BookMuted!30,3/FigureRed!70,4/FigureRed!70}{
    \fill[\col] (\x-0.4,1.6) rectangle (\x+0.4,2.4);
    \draw[BookInk,line width=0.5pt] (\x-0.4,1.6) rectangle (\x+0.4,2.4);
    \node[below] at (\x,1.55) {$\x$};
  }
  \node[anchor=east] at (-0.6,2.0) {像素};
  % 核权重
  \node[anchor=east] at (-0.6,0.75) {核权重};
  \foreach \x/\k in {0/0.36,1/0.84,2/1.00,3/0.84,4/0.36}{
    \node at (\x,0.75) {$\k$};
  }
  % 像素权重
  \node[anchor=east] at (-0.6,0.15) {$w_i$};
  \foreach \x/\w in {0/0.786,1/0.786,2/0.786,3/1.304,4/1.304}{
    \node at (\x,0.15) {$\w$};
  }
  % 窗口中心与移动
  \draw[BookTeal,line width=1.1pt] (2,2.55) -- (2,3.05);
  \node[above] at (2,3.05) {$y_0=2$};
  \draw[BookBlue,line width=1.1pt] (2.31,2.55) -- (2.31,3.05);
  \draw[arr] (2.02,2.8) -- (2.29,2.8);
  \node[anchor=west] at (2.4,3.2) {$y_1\approx2.31$};
  % 带宽范围
  \draw[BookTeal,line width=0.7pt,{Bar[width=2mm]}-{Bar[width=2mm]}] (-0.5,-0.45) -- (4.5,-0.45);
  \node[below] at (2,-0.5) {带宽$h=2.5$覆盖的窗口};
\end{tikzpicture}
